% !TeX root = ../main.tex
\section{Introduction}
\label{sec:intro}

A guitar part on a record is the product of the notes and of the equipment
that shaped them: the guitar and its pickups, often a compressor or a drive
pedal, an amplifier whose preamplifier, tone stack and power stage distort and
filter the signal, a speaker cabinet recorded by a microphone, modulation,
delay and reverberation, and the mixing and mastering applied afterwards.
Guitarists call the combined result the \emph{tone}. A player learning a song
wants the notes, which tablature transcription addresses~\cite{deangelo2026tabreview},
and also the tone, which today is found by ear, by reading interviews about
the artist's equipment, or by buying a product preset made for that song.

Estimating the tone automatically is an inverse problem with several layers
of difficulty. The processing is nonlinear (distortion), time-varying
(modulation, compression) and long-memory (delay, reverberation). The
mapping from settings to sound is not one-to-one: amplifier gain, a drive
pedal and a compressor can produce similar sounds, and so can different
orders of the same effects. In a commercial song the guitar is mixed with
other instruments, so it must first be separated, and separation leaves
artefacts that look like effects. The dry signal that entered the chain is
never available, which rules out the standard way of modelling an amplifier
from paired input and output. And the target is not to identify the original
equipment but to approximate its sound with the equipment the player has,
here the effects chain of a real-time application.

This review is written for a concrete system. Guitar Virtual Amp is an
open-source, real-time effects processor with a mono chain of 18 effects
(amplifier with three voicings and a tone stack, distortion, compressor,
noise gate, equaliser, chorus, flanger, phaser, tremolo, wah, octaver, pitch
shifter, delay, reverb, low-pass filter, looper, an impulse-response cabinet
and a Neural Amp Modeler capture player). It already separates the guitar from
a song with a band-split transformer, splits it by stereo position, and
transcribes the notes and the beats~\cite{deangelo2026tabreview}. The question
is what to add so that it can propose a preset that sounds like the guitar in
a recording the user owns.

Fig.~\ref{fig:families} groups the relevant work into five families. Each
family solves part of the problem under assumptions that the commercial case
breaks, usually the availability of a dry signal, synthetic training data, or
an isolated guitar.

\begin{figure}[t]
\centering
\begin{tikzpicture}[
  node distance=1.6mm and 2mm,
  box/.style={draw, rounded corners=1pt, font=\scriptsize, align=center,
              minimum height=5mm, inner sep=2pt, fill=white, text width=17mm},
  wide/.style={box, text width=27mm},
  arr/.style={-{Stealth[length=1.6mm]}, semithick}]
  \node[box] (stem) {Guitar stem\\(no DI)};
  \node[wide, right=6mm of stem, yshift=15mm] (rec) {(a) Recognition and\\parameter regression};
  \node[wide, below=of rec] (abs) {(b) Analysis by synthesis\\(search over a chain)};
  \node[wide, below=of abs] (unp) {(c) Unpaired / blind\\neural modelling};
  \node[wide, below=of unp] (rem) {(d) Effect removal,\\then modelling};
  \node[wide, below=of rem] (ret) {(e) Retrieval of a\\preset or capture};
  \node[box, right=6mm of abs, yshift=-7mm] (out) {Preset, parameters\\or neural model};
  \foreach \n in {rec,abs,unp,rem,ret} {
    \draw[arr] (stem.east) -- (\n.west);
    \draw[arr] (\n.east) -- (out.west);
  }
\end{tikzpicture}
\caption{Families of methods that estimate a tone from processed audio. (a)
predicts effect types and settings with a network; (b) renders candidate
settings and compares them with the target; (c) learns a neural model of the
device from processed audio alone; (d) recovers a dry signal and then fits a
model on the recovered pair; (e) returns the closest item of a library.}
\label{fig:families}
\end{figure}

We cover the task and its metrics (Section~\ref{sec:task}), datasets and
licences (Section~\ref{sec:datasets}), effect recognition and parameter
estimation (Section~\ref{sec:recognition}), amplifier and effect modelling
(Section~\ref{sec:modelling}), methods that work without the dry signal
(Section~\ref{sec:nodi}), style transfer, effect representations and the
lessons of synthesizer sound matching (Section~\ref{sec:matching}), products
and open-source tools (Section~\ref{sec:products}), a comparison with failure
modes and gaps (Section~\ref{sec:comparison}), and an approach for the
application with an evaluation protocol (Section~\ref{sec:application}).
Numbers are quoted as published; results on different data are not directly
comparable, and we say so where it matters. Statements about products come
from their manuals and pages; vendors publish no accuracy figures for
recording-based tone features.
